\documentclass[11pt,a4paper]{article}
\usepackage[margin=25mm]{geometry}
\usepackage[T1]{fontenc}
\usepackage{lmodern,amsmath,amssymb,amsthm}
\usepackage{microtype}
\usepackage{needspace}
\usepackage[all]{xy}
\usepackage{fancyhdr}
\usepackage[hidelinks]{hyperref}
\hypersetup{
pdftitle={Jardim: exercise-session notes},
pdfsubject={Rao modules, common sheaves,
and a sheaf primer},
pdfauthor={}}
\pagestyle{fancy}
\fancyhf{}
\fancyhead[L]{\small Jardim: exercise notes}
\fancyhead[R]{\small Salvador, 2026}
\fancyfoot[C]{\thepage}
\setlength{\headheight}{14pt}
\setlength{\parindent}{0pt}
\setlength{\parskip}{5pt}
\setlength{\emergencystretch}{2em}
\newtheorem{proposition}{Proposition}[section]
\newtheorem{lemma}[proposition]{Lemma}
\theoremstyle{remark}
\newtheorem{remark}[proposition]{Remark}

% Build from this directory with:
% latexmk -pdf -outdir=build \
%   jardim-exercises.tex

\begin{document}
\thispagestyle{plain}
\begin{center}
{\LARGE Curves, common sheaves and
Rao modules}\par
\vspace{6pt}
{\large Notes from Marcos Jardim's
exercise sessions}\par
\vspace{5pt}
IV Brazilian Workshop on Lefschetz
Properties\par
UFBA, Salvador --- September 28--October 1,
2026
\end{center}
\vspace{4pt}

These notes combine Larissa's session
notes and the blackboard photographs.
The common-plane argument was explained
by Thiago Fassarella; the sheaf-Ext
criteria were explained by Charles.
The explanations between the calculations
are added to make the arguments readable
without the blackboard. Appendix~A recalls
the sheaf terminology and explains Ext.
The monad in Section~4 remains a proposed
construction with extra conditions to check.

The question running through the exercises
is this: \emph{when can two curves be
obtained from sections of the same sheaf?}
We work on projective three-space:
$$
X=\mathbb{P}^3_{\mathbb{C}},\qquad
R=\mathbb{C}[x_0,x_1,x_2,x_3],\qquad
h=c_1(\mathcal{O}_X(1)).
$$
Write $\mathcal{O}(a)=\mathcal{O}_X(a)$,
$F(a)=F\otimes\mathcal{O}(a)$ and
$$
H^i_*(F)=\bigoplus_{t\in\mathbb{Z}}
H^i(X,F(t)).
$$
The star in $H^i_*(F)$ means that we keep
the cohomology of \emph{every twist},
with the integer $t$ recording the degree.
For a graded module $N$, our shift
convention is $N[e]_t=N_{t+e}$.
Sheaf duals are denoted by $F^*$ and
duals of vector spaces by $V^\vee$.

\section{From a section to a curve}

Let $E$ be a rank-two vector bundle
with $\det E=\mathcal{O}(c)$.
The integer $c$ records its first Chern
class: $c_1(E)=ch$.
A section of $E(a)$ is initially a map
$\mathcal{O}\to E(a)$: it sends the
constant section $1$ to the chosen
section $s$. Tensoring with
$\mathcal{O}(-a)$ gives the equivalent
map $s:\mathcal{O}(-a)\to E$.
This is why the minus sign appears.

Suppose the zero scheme of $s$ is a
curve $C$, of pure codimension two.
Then
\begin{equation}
0\longrightarrow\mathcal{O}(-a)
\xrightarrow{s}E
\longrightarrow\mathcal{I}_C(c+a)
\longrightarrow0.
\label{eq-section}
\end{equation}
There are two things to understand here.
First, where $E$ is trivial, the section
is a pair of functions $s=(f,g)$ and
$C=V(f,g)$. The map
$$
(u,v)\longmapsto fv-gu
$$
has image the ideal $(f,g)$.
Since $f,g$ form a regular sequence,
its kernel is generated by $(f,g)$.
Thus quotienting by the section really
does give the ideal of its zero curve.
This also shows that $C$ is locally a
complete intersection.

Second, the twist of that ideal is
forced by the determinant. If it were
$\mathcal{I}_C(b)$, then away from $C$
the determinant identity would be
$$
\mathcal{O}(c)
\cong\mathcal{O}(-a)\otimes\mathcal{O}(b).
$$
Hence $b=c+a$. The ideal sheaf is
torsion-free, so the inclusion of the
section has \emph{saturated} image.
This does not say that its quotient
is a line bundle along $C$.

The starting question is whether two
curves can occur in sequences of the
form (1) with the same middle sheaf.
It matters whether this middle sheaf
must be a vector bundle or may instead
be reflexive. The distinction will be
essential for the line and the conic.

\Needspace{7\baselineskip}
\textbf{What the Rao module measures.}
For a locally Cohen--Macaulay curve
$C\subset X$, its \emph{Rao module} is
$$
M_C=H^1_*(\mathcal{I}_C).
$$
It is a finite-length graded $R$-module.
Multiplication by $f\in R_s$ sends its
degree-$t$ part to its degree-$(t+s)$
part through the map on cohomology
induced by multiplication by $f$.
The exact sequence
$$
0\longrightarrow\mathcal{I}_C(t)
\longrightarrow\mathcal{O}(t)
\longrightarrow\mathcal{O}_C(t)
\longrightarrow0
$$
explains its meaning. For $t\geq0$,
$$
(M_C)_t\cong\operatorname{coker}
\bigl(R_t\longrightarrow
H^0(C,\mathcal{O}_C(t))\bigr).
$$
Thus it measures sections on $C$ that
do not extend to homogeneous polynomials
on the ambient projective space.
Its module structure remembers how these
classes change under multiplication by
polynomials, not just their dimensions.
The curve is arithmetically
Cohen--Macaulay exactly when $M_C=0$.

\begin{proposition}
If $C_1,C_2$ occur in (1) with a common
$E$ and twists $a_1,a_2$, then
$$
M_{C_1}[c+a_1]\cong M_{C_2}[c+a_2],
\qquad
M_{C_1}\cong M_{C_2}[a_2-a_1].
$$
\end{proposition}
\begin{proof}
For every integer $m$,
$H^1(\mathcal{O}(m))=
H^2(\mathcal{O}(m))=0$.
Twisting (1) by $t$ therefore gives
$$
0\longrightarrow H^1(E(t))
\longrightarrow
H^1(\mathcal{I}_{C_i}(t+c+a_i))
\longrightarrow0.
$$
The zero on the left is
$H^1(\mathcal{O}(t-a_i))$;
the zero on the right is
$H^2(\mathcal{O}(t-a_i))$.
Consequently
$$
H^1(E(t))\cong
H^1(\mathcal{I}_{C_i}(t+c+a_i)).
$$
These isomorphisms commute with
multiplication by homogeneous
polynomials. Both curves therefore give
the same graded module $H^1_*(E)$,
with different choices of where degree
zero sits.
\end{proof}

\Needspace{7\baselineskip}
\textbf{What Rao's theorem says.}
Rao's classification says that two
locally Cohen--Macaulay curves in
$\mathbb{P}^3$ belong to the same
\emph{even liaison class} if and only
if their Rao modules are isomorphic
up to shift \cite{rao}.
Even liaison means a chain of an even
number of complete-intersection links.
This classification does not assert
the existence of one common rank-two
vector bundle in (1). A common bundle
implies equality of shifted Rao modules;
Rao converts that equality into a chain
of links. The reverse passage to bundles
allows changing rank by adding sums of
line bundles, as well as twisting.

\textbf{Higher-rank version.}
If the bundle has rank $r$, a single
section generally leaves a quotient of
rank $r-1$. To obtain an ideal sheaf,
which has rank one, we instead use
$r-1$ sections and sequences
$$
0\longrightarrow
\mathcal{O}(-a_i)^{\oplus(r-1)}
\longrightarrow E
\longrightarrow\mathcal{I}_{C_i}(b_i)
\longrightarrow0,
\qquad b_i=c+(r-1)a_i.
$$
When the degeneracy locus has
codimension two, it is defined locally
by the maximal minors of the matrix
of sections. The same proof gives
$$
M_{C_1}[b_1]\cong M_{C_2}[b_2].
$$
The classical stable correspondence
allows sums of line bundles of different
degrees and stabilization by such sums.
It does not prescribe this particular
rank and these particular kernels.

\section{Common quotients and extensions}

The next diagram is a way to compare
two presentations without discarding
the maps between the sheaves.
Suppose $E$ has two
saturated line-sheaf inclusions
$A_i=\mathcal{O}(-a_i)\to E$, with
torsion-free quotients $Q_i$:
\begin{equation}
0\longrightarrow A_i\xrightarrow{s_i}E
\longrightarrow Q_i\longrightarrow0.
\end{equation}
Then
$$
H^1_*(Q_1)\cong H^1_*(E)
\cong H^1_*(Q_2).
$$
There is no shift here unless the
quotients themselves are retwisted.
The converse asks whether equality of
these cohomology modules can produce
a common middle sheaf with the
specified kernels.

\Needspace{13\baselineskip}
\textbf{Quotienting by both sections.}
We have $Q_1=E/s_1(A_1)$ and
$Q_2=E/s_2(A_2)$. If we quotient by
both images, in either order, we get
$$
T=E/(s_1(A_1)+s_2(A_2)).
$$
If the sections are independent over
the function field, neither disappears
when we quotient by the other. This
gives the diagram
$$
\xymatrix{
&A_1\ar@{=}[r]\ar[d]^{s_1}&A_1\ar[d]\\
A_2\ar@{=}[d]\ar[r]^{s_2}
&E\ar[r]\ar[d]&Q_2\ar[d]\\
A_2\ar[r]&Q_1\ar[r]&T
}
$$
has four short exact rows and columns,
with zeros at their ends understood.
In particular, the bottom row and the
right column say
$$
Q_1/A_2\cong T\cong Q_2/A_1.
$$
The two inclusions into the quotients
are the original sections followed by
the quotient maps.
A nonzero map from a line bundle to a
torsion-free sheaf is injective.
In particular, if $a_1<a_2$, the map
$A_1\to Q_2$ cannot vanish: vanishing
would factor $s_1$ through a map
$$
A_1\longrightarrow A_2,
\qquad
\operatorname{Hom}(A_1,A_2)
=H^0(\mathcal{O}(a_1-a_2))=0.
$$
At $a_1=a_2$, proportional sections
must be excluded.

\Needspace{6\baselineskip}
\textbf{Reading the diagram backwards.}
Conversely, given surjections $Q_i\to T$
with kernels $A_2,A_1$, respectively,
the fibre product $E=Q_1\times_TQ_2$
has both sequences (2). Its local
sections are pairs $(q_1,q_2)$ whose
images in $T$ agree. Projection onto
$Q_1$ has kernel $A_1$, and projection
onto $Q_2$ has kernel $A_2$.
Whether this $E$ is locally free or
reflexive is a separate local question.
The analogous graded-module construction
uses injections $R[-a_i]\to M$ and
their cokernels; sheafification gives
the same diagram on projective space.

\Needspace{8\baselineskip}
\textbf{How this becomes ordinary linkage.}
In rank two, write
$Q_i=\mathcal{I}_{C_i}(b_i)$.
The determinant $s_1\wedge s_2$ defines
a surface $S$ containing both curves.
Here assume the sections are generically
independent and, for simplicity, that
$S$ is smooth. The common quotient is
$$
\mathcal{I}_{C_1/S}(b_1)
\cong T\cong\mathcal{I}_{C_2/S}(b_2),
$$
where we view sheaves on $S$ as sheaves
on projective space by extension by zero.
On a smooth surface the ideal of a
curve is the line bundle
$\mathcal{O}_S(-C_i)$. Thus, if $H_S$
is the hyperplane class of $S$,
$$
C_2\sim C_1+(b_2-b_1)H_S.
$$
This is an elementary \emph{biliaison}.
To see the two actual links, choose a
general residual curve $D$ such that
$C_1+D$ is a sufficiently high-degree
hypersurface section of $S$.
The displayed equivalence says that
$C_2+D$ is also a hypersurface section
of $S$. Hence both unions are complete
intersections in $\mathbb{P}^3$:
first link $C_1$ to $D$, then $D$ to
$C_2$. For example, if $S=V(F)$ and
$C_1\cup D=V(F,G)$, the first link is
$$
I_D=(F,G):I_{C_1}.
$$
This is the residual-ideal construction.
Hartshorne's Proposition~2.5 treats the
sheaf argument using generalized divisors,
so the surface need not be smooth
\cite{rao}.

\Needspace{7\baselineskip}
\textbf{Extension spaces.}
Sequences (2), with end terms fixed,
are classified by
$\operatorname{Ext}^1(Q_i,A_i)$.
An element of this group specifies how
to assemble $A_i$ and $Q_i$ into a
middle sheaf $E$. The zero element
gives the direct sum; other elements
give nonsplit extensions.
Serre duality gives
$$
\operatorname{Ext}^1
(Q_i,\mathcal{O}(-a_i))
\cong H^2(Q_i(a_i-4))^\vee.
$$
Thus an isomorphism of $H^1_*(Q_i)$
does not directly identify the
extension spaces. Even an isomorphism
of these spaces would not by itself
identify the middle sheaves.
Section~4 gives a concrete obstruction
to the proposed converse.

\section{A line and a conic}

Let $L$ be a line and $Q$ a smooth
plane conic in $X$.
The line is cut out by two linear
equations; the conic by one linear and
one quadratic equation. Each pair has
one relation, giving the resolutions
$$
\begin{aligned}
0&\longrightarrow\mathcal{O}(-2)
\longrightarrow\mathcal{O}(-1)^{\oplus2}
\longrightarrow\mathcal{I}_L
\longrightarrow0,\\
0&\longrightarrow\mathcal{O}(-3)
\longrightarrow
\mathcal{O}(-1)\oplus\mathcal{O}(-2)
\longrightarrow\mathcal{I}_Q
\longrightarrow0.
\end{aligned}
$$
For instance, if $I_L=(f,g)$, its relation
is $-gf+fg=0$. The coefficients have
degree one, so the relation between the
two degree-one generators has total
degree two: this explains the
$\mathcal{O}(-2)$ on the left.
The conic has a relation of total
degree three for the same reason.
Taking cohomology and using the
vanishing of $H^1$ and $H^2$ for line
bundles gives $M_L=M_Q=0$.
Thus Rao's theorem puts them in the same
even liaison class. We now ask the
stronger question about a common sheaf.

\Needspace{7\baselineskip}
\textbf{Computing the extensions.}
Fixing $a$ and $c$ fixes the two end
terms in (1). We can first ask how
many extension classes there are.
For $C=L$ or $Q$, Appendix~A explains
each step of
$$
\begin{aligned}
\operatorname{Ext}^1
(\mathcal{I}_C(c+a),\mathcal{O}(-a))
&\cong H^2(\mathcal{I}_C(c+2a-4))^\vee\\
&\cong H^1(\mathcal{O}_C(c+2a-4))^\vee.
\end{aligned}
$$
This reduces Ext on three-space to
line-bundle cohomology on a curve.
Both curves are isomorphic to
$\mathbb{P}^1$, but their hyperplane
bundles are different: a hyperplane
meets $L$ in one point and $Q$ in two.
Thus $\mathcal{O}_L(m)$ corresponds to
$\mathcal{O}_{\mathbb{P}^1}(m)$, whereas
$\mathcal{O}_Q(m)$ corresponds to
$\mathcal{O}_{\mathbb{P}^1}(2m)$.
The usual formula on $\mathbb{P}^1$ gives
$$
\begin{aligned}
h^1(\mathcal{O}_L(m))&=\max\{0,-m-1\},\\
h^1(\mathcal{O}_Q(m))&=\max\{0,-2m-1\}.
\end{aligned}
$$
Thus the two extension spaces have
dimensions
$$
\max\{0,3-c-2a_1\},\qquad
\max\{0,7-2c-4a_2\},
$$
respectively. A nonzero class gives
a nonsplit extension, not necessarily
a locally free middle sheaf.

\Needspace{8\baselineskip}
\begin{proposition}
A line and a smooth conic cannot be
zero schemes of sections of twists
of one rank-two vector bundle.
\end{proposition}
\begin{proof}
If such a bundle existed, adjunction
would give
$$
\omega_C\cong\mathcal{O}_C(c+2a-4).
$$
The $-4$ comes from
$\omega_{\mathbb{P}^3}=\mathcal{O}(-4)$,
and $c+2a$ from $\det(E(a))$.
Both curves have canonical bundle
$\mathcal{O}_{\mathbb{P}^1}(-2)$.
For the line we must therefore have
$c+2a_1-4=-2$.
For the conic the twist is doubled on
$\mathbb{P}^1$, so we must have
$2(c+2a_2-4)=-2$.
Equivalently,
$$
c+2a_1=2,\qquad c+2a_2=3.
$$
Subtracting gives $2(a_2-a_1)=1$,
which is impossible.
\end{proof}

\Needspace{8\baselineskip}
\textbf{The common plane
(Thiago Fassarella's argument).}
We can still ask whether a rank-two
reflexive sheaf might work: such a
sheaf is allowed to fail to be a
bundle at finitely many points.
Suppose there is a rank-two
\emph{reflexive} sheaf $E$ with
$\det E=\mathcal{O}(c)$ and sequences
\begin{equation}
\begin{aligned}
0&\longrightarrow\mathcal{O}(-a_1)
\xrightarrow{s_1}E
\longrightarrow\mathcal{I}_L(c+a_1)
\longrightarrow0,\\
0&\longrightarrow\mathcal{O}(-a_2)
\xrightarrow{s_2}E
\longrightarrow\mathcal{I}_Q(c+a_2)
\longrightarrow0.
\end{aligned}
\end{equation}
We will first find the degree of the
determinant equation, then show that
it vanishes on both curves.

\textbf{Step 1: compare degrees through
Chern classes.}
The total Chern class of a line bundle
$\mathcal{O}(t)$ is $1+th$.
Total Chern classes multiply in exact
sequences, so the two resolutions give
$$
\begin{aligned}
c(\mathcal{I}_L(t))
&=\frac{(1+(t-1)h)^2}{1+(t-2)h},\\
c(\mathcal{I}_Q(t))
&=\frac{(1+(t-1)h)(1+(t-2)h)}
{1+(t-3)h}.
\end{aligned}
$$
Expand these fractions and discard
powers $h^4$ and higher, since they
vanish on $\mathbb{P}^3$. The result is
$$
\begin{aligned}
c(\mathcal{I}_L(t))
&=1+th+h^2+(2-t)h^3,\\
c(\mathcal{I}_Q(t))
&=1+th+2h^2+(6-2t)h^3.
\end{aligned}
$$
Multiplying by $1-a_i h$ in (3) and
comparing the coefficients of $h^2$
gives
$$
1-a_1(c+a_1)=2-a_2(c+a_2),
$$
or equivalently
\begin{equation}
\boxed{1=(a_2-a_1)(c+a_1+a_2).}
\end{equation}
The $1$ on the left is the difference
between the degrees of the conic and
the line. Both presentations must give
the same $c_2(E)$; that is the only
input in this equation.

\Needspace{7\baselineskip}
\textbf{Step 2: take the determinant
of the two sections.}
The determinant of a reflexive sheaf
is $(\bigwedge^2E)^{**}$.
Consequently the two sections give
$$
s_1\wedge s_2\in
H^0(\mathcal{O}(d)),\qquad
d=c+a_1+a_2.
$$
To track the twist, the wedge is a map
from $\mathcal{O}(-a_1-a_2)$ to
$\det E=\mathcal{O}(c)$, so it is a
homogeneous polynomial of degree $d$.
Locally, where $E$ is a bundle, it is
just the determinant of the two columns
$s_1,s_2$.
This section is nonzero. Otherwise
the two saturated images would be
the same line over the function field,
and hence the same subsheaf of $E$.
This would force $a_1=a_2$, contrary
to (4). Therefore $d\geq0$, and (4)
forces
$$
d=1,\qquad a_2-a_1=1.
$$
The alternative factorization
$1=(-1)(-1)$ is ruled out because a
negative-degree line bundle has no
nonzero global sections. Thus the
determinant is a nonzero linear
polynomial, and its zero divisor is
a plane:
$$
\boxed{H=Z(s_1\wedge s_2),\qquad
L\cup Q\subset H.}
$$
To see the containment, compose $s_2$
with the quotient to
$\mathcal{I}_L(c+a_1)$.
This gives a section of
$\mathcal{I}_L(d)$.
The other composition gives a section
of $\mathcal{I}_Q(d)$.
Both are the determinant section,
up to sign, so its linear equation
belongs to both ideals.
On the locally free locus this is
just the determinant of two column
vectors, which vanishes wherever
either section vanishes.

\Needspace{7\baselineskip}
\textbf{Step 3: record the remaining
constraints.}
Substituting $a_2=a_1+1$ into
$c+a_1+a_2=1$ gives $c=-2a_1$.
Using the line's sequence, we also get,
writing
$c_i(E)$ for its integer coefficient,
$$
c=-2a_1,\qquad c_2(E)=1+a_1^2,
\qquad c_3(E)=2.
$$
For the last equality, the coefficient
of $h^3$ in the line's sequence is
$2-(c+a_1)-a_1=2-c-2a_1=2$.
The third Chern class again rules out
local freeness. These are necessary
conditions for (3); they do not yet
construct the common sheaf.

\Needspace{8\baselineskip}
\textbf{The trial used in the exercises.}
For $a_1=1,a_2=2,c=-2$,
$$
\begin{aligned}
0&\longrightarrow\mathcal{O}(-1)
\longrightarrow E
\longrightarrow\mathcal{I}_L(-1)
\longrightarrow0,\\
0&\longrightarrow\mathcal{O}(-2)
\longrightarrow E
\longrightarrow\mathcal{I}_Q
\longrightarrow0.
\end{aligned}
$$
Here $c_2(E)=2$ and both extension
spaces have dimension three: they
reduce to $H^1(\mathcal{O}_L(-4))$
and $H^1(\mathcal{O}_Q(-2))$, which
both identify, via the parametrizations,
with
$H^1(\mathcal{O}_{\mathbb{P}^1}(-4))$.
The induced maps are
$\mathcal{O}(-1)\to\mathcal{I}_Q$
and $\mathcal{O}(-2)\to\mathcal{I}_L(-1)$.
They are multiplication by the equation
of the common plane $H$.
Quotienting by this equation means
working inside $H$. There, a line has
ideal $\mathcal{O}_H(-1)$ and a conic
has ideal $\mathcal{O}_H(-2)$.
The extra twist on the line makes
the two quotients agree.
For $i:H\hookrightarrow X$, this says
$$
T\cong i_*\mathcal{I}_{L/H}(-1)
\cong i_*\mathcal{O}_H(-2)
\cong i_*\mathcal{I}_{Q/H}.
$$
This is the rank-two instance of the
quotient diagram in Section~2.

More generally,
$H^0(\mathcal{I}_C(k))\ne0$, for $k>0$,
means that $C$ lies in a degree-$k$
hypersurface: a nonzero section is
a homogeneous polynomial vanishing
on $C$.
For a coplanar line and conic with
equations $L=V(\ell,\ell')$ and
$Q=V(\ell,q)$, their union is
$V(\ell,\ell'q)$ and has degree three.

\Needspace{8\baselineskip}
\textbf{Dualizing and locating
singularities.}
For $C=L,Q$ and $t=c+a$, dualizing
the corresponding sequence gives
$$
\begin{aligned}
0&\longrightarrow\mathcal{O}(-t)
\longrightarrow E^*
\longrightarrow\mathcal{O}(a)\\
&\longrightarrow\mathcal{E}xt^1
(\mathcal{I}_C(t),\mathcal{O})
\longrightarrow\mathcal{E}xt^1
(E,\mathcal{O})\longrightarrow0.
\end{aligned}
$$
The curve resolutions identify
$$
\begin{aligned}
\mathcal{E}xt^1
(\mathcal{I}_L(t),\mathcal{O})
&\cong\mathcal{O}_L(2-t),\\
\mathcal{E}xt^1
(\mathcal{I}_Q(t),\mathcal{O})
&\cong\mathcal{O}_Q(3-t).
\end{aligned}
$$
Hence the support of
$\mathcal{E}xt^1(E,\mathcal{O})$
is contained in $L\cap Q$: this sheaf
is a quotient of a sheaf supported
on $L$, and also a quotient of a sheaf
supported on $Q$. For reflexive $E$,
its support is exactly the set where
$E$ is not locally free.
So any singular points of the common
sheaf must lie in the intersection
of the two curves. Appendix~A explains
why sheaf Ext detects this failure.

\section{Higher rank: an obstruction
and a monad}

\textbf{Equal cohomology is insufficient.}
The problem with trying to reconstruct
a sheaf from $H^1_*$ is that line-bundle
summands are invisible to this module.
They have no $H^1$ in any twist, but
they do change Chern classes.
A concrete version of the session's
Chern-class example uses the cotangent
bundle $F=\Omega^1_X$ and
$$
Q_1=F\oplus\mathcal{O}^{\oplus2},
\qquad Q_2=F\oplus\mathcal{O}(1)
\oplus\mathcal{O}(-1).
$$
The Euler sequence is
$$
0\longrightarrow\Omega^1_X
\longrightarrow\mathcal{O}(-1)^{\oplus4}
\longrightarrow\mathcal{O}
\longrightarrow0.
$$
Thus $c(F)=(1-h)^4$ modulo $h^4$, and
$$
\begin{aligned}
c(Q_1)&=1-4h+6h^2-4h^3,\\
c(Q_2)&=c(F)(1+h)(1-h)
=1-4h+5h^2.
\end{aligned}
$$
Yet
$$
H^1_*(Q_1)\cong H^1_*(F)
\cong H^1_*(Q_2),
$$
since line bundles have no $H^1_*$.
If both appeared in (2) with a common
$E$, then
$$
(1-a_1h)c(Q_1)=(1-a_2h)c(Q_2).
$$
The coefficient of $h$ forces
$a_1=a_2$, and the coefficient of
$h^2$ then gives a contradiction.
This disproves the converse for
these fixed quotients and
line-bundle kernels.

\Needspace{8\baselineskip}
\textbf{Starting from a module.}
We can also run the construction in
the other direction: start with a
module and seek a sheaf having that
module as $H^1_*$.
Let $M$ be a nonzero finite-length
graded $R$-module, with a minimal
free presentation
$$
L_1\xrightarrow{\beta}L_0
\longrightarrow M\longrightarrow0,
\qquad L_0=
\bigoplus_iR[-a_i]^{\oplus m_i}.
$$
The tilde takes a graded module to
its associated sheaf on
$\operatorname{Proj}R$.
A finite-length graded module is
supported at the irrelevant ideal
$(x_0,x_1,x_2,x_3)$, which is not a
point of projective space.
Consequently $\widetilde M=0$, so
$\widetilde\beta:\widetilde L_1
\to\widetilde L_0$ is surjective.
Its kernel $K$ is a vector bundle:
locally a surjection onto a free
module splits. We have
$$
0\longrightarrow K
\longrightarrow\widetilde L_1
\longrightarrow\widetilde L_0
\longrightarrow0.
$$
The long exact cohomology sequence,
in all twists, now gives
$$
L_1\longrightarrow L_0
\longrightarrow H^1_*(K)
\longrightarrow0.
$$
Thus $H^1_*(K)\cong M$.
This first step already produces a
bundle with the desired cohomology;
its rank and Chern classes may not
be the ones we want.

\Needspace{8\baselineskip}
\textbf{What the monad would add.}
The further construction discussed
in the session asks for an injection
$$
\alpha:\widetilde L_0^*
\longrightarrow\widetilde L_1,
\qquad\widetilde\beta\alpha=0.
$$
The equation $\widetilde\beta\alpha=0$
means that the image of $\alpha$ lies
inside $K$. The resulting complex
$$
\widetilde L_0^*
\xrightarrow{\alpha}\widetilde L_1
\xrightarrow{\widetilde\beta}
\widetilde L_0
$$
is called a \emph{monad}: it is exact
at the two ends, and its middle
cohomology is
$$
E=\ker\widetilde\beta/
\operatorname{im}\alpha.
$$
Equivalently,
$$
0\longrightarrow\widetilde L_0^*
\longrightarrow K\longrightarrow E
\longrightarrow0.
$$
Because the left term is a sum of
line bundles, it has no $H^1_*$ or
$H^2_*$. Thus
$$
H^1_*(E)\cong H^1_*(K)
\cong\operatorname{coker}(L_1\to L_0)
\cong M.
$$
The needed map $\alpha$ is
additional data: dualizing $\beta$
lands in $\widetilde L_1^*$, not
automatically in $\widetilde L_1$.
For $E$ to be locally free, the image
of $\alpha$ must be a subbundle of $K$.
The ranks and first Chern classes are
$$
\begin{aligned}
\operatorname{rk}E
&=\operatorname{rk}L_1
-2\operatorname{rk}L_0,\\
c_1(E)&=c_1(\widetilde L_1).
\end{aligned}
$$
In the second equality the first Chern
classes of $\widetilde L_0$ and its
dual cancel. Hence the intended
condition $c_1(E)=0$ must be checked
on $\widetilde L_1$.
Constructing this monad and suitable
saturated sections would then produce
curves whose Rao module is $M$, up to
shift.

\section{The Lefschetz question}

For a finite-length graded $R$-module
$N$, the \emph{weak Lefschetz property}
asks for $\ell\in R_1$ such that
$$
\times\ell:N_t\longrightarrow N_{t+1}
$$
has maximal rank for every $t$.
The \emph{strong Lefschetz property}
asks that the same $\ell$ make
$$
\times\ell^s:N_t\longrightarrow N_{t+s}
$$
have maximal rank for every $t$ and
$s\geq1$.
Maximal rank means that the rank is
the smaller of the dimensions of the
source and target. Thus the map must
be injective when the source is
smaller, and surjective when the
target is smaller. One choice of
$\ell$ must work in all degrees;
for the strong property it must also
work for every power.
Jardim's question is which curves
have Rao modules with these properties.
Both properties survive graded
isomorphisms and shifts, hence are
invariants of the even liaison class.
For ACM curves the Rao module is zero,
so both hold vacuously.

\clearpage
\appendix
\section{Sheaf primer}
\label{appendix-primer}

Here sheaves are coherent. Statements
about rank and torsion use an integral
variety; the Ext tests at the end use
a smooth projective threefold.

\Needspace{7\baselineskip}
\textbf{Coherent, locally free, bundle.}
A coherent sheaf is locally described
by finitely many generators and
relations. On an affine open set
$\operatorname{Spec}R$, it corresponds
to a finite $R$-module.
It is \emph{locally free of rank $r$}
if every point has a neighbourhood $U$
with
$$
F|_U\cong\mathcal{O}_U^{\oplus r}.
$$
This means that its sections locally
have a basis of $r$ elements. It is
exactly the sheaf of sections of an
algebraic vector bundle; a basis of
sections is a frame of the bundle.
Rank-one bundles are line bundles,
also called invertible sheaves.
The rank of an arbitrary coherent
sheaf is its dimension over the
function field at the generic point.
It describes the sheaf on a dense
open set, and does not guarantee a
basis near every point.

\Needspace{8\baselineskip}
\textbf{Torsion concerns multiplication
by scalars.}
Over a domain $R$, an element $m\in M$
is torsion if $fm=0$ for some nonzero
$f\in R$. The module is torsion-free
if this forces $m=0$.
For example, take
$$
R=\mathbb{C}[x],\qquad M=R/(x).
$$
The element $\overline 1\in M$ is
nonzero, but $x\overline 1=0$.
The scalar $x$ is nonzero in $R$;
its class in the quotient is zero.
That is precisely how a nonzero
scalar kills a nonzero module element.
There are no nonzero nilpotents in
$M\cong\mathbb{C}$. Torsion concerns
the $R$-action on $M$, whereas
nilpotence concerns powers of an
element inside a ring.

For sheaves, apply the module definition
to each stalk. If $Z$ is a proper closed
subscheme of an integral variety, then
$$
0\longrightarrow\mathcal{I}_Z
\longrightarrow\mathcal{O}_X
\longrightarrow\mathcal{O}_Z
\longrightarrow0
$$
has a torsion-free ideal sheaf: it is
a submodule of a domain at each stalk.
Its quotient is torsion as a sheaf on
$X$, since local equations of $Z$
annihilate it.

\Needspace{9\baselineskip}
\textbf{The canonical map to the
 double dual.}
The dual is
$$
F^*=\mathcal{H}om(F,\mathcal{O}_X).
$$
For modules, an element of $M^*$ is
an $R$-linear map $\varphi:M\to R$.
Given $m\in M$, we can evaluate every
such map at $m$; this gives
$$
\begin{aligned}
\operatorname{ev}:M&\longrightarrow M^{**},\\
\operatorname{ev}(m)(\varphi)&=\varphi(m).
\end{aligned}
$$
This is canonical because it makes no
choice of basis. For a finite-dimensional
vector space it is an isomorphism.
For a module over a ring it may fail
to be injective or surjective.
A sheaf is \emph{reflexive} precisely
when this evaluation map is an
isomorphism. For coherent sheaves on
an integral Noetherian scheme,
torsion-freeness is equivalent to its
being injective \cite{stacks}.

The implication from reflexive to
torsion-free has a short proof.
Suppose $fm=0$ with $f\ne0$.
For every $\varphi\in M^*$,
$$
f\varphi(m)=\varphi(fm)=0.
$$
The ring is a domain, so
$\varphi(m)=0$ for every $\varphi$.
Thus $\operatorname{ev}(m)=0$, and
injectivity of evaluation gives $m=0$.
In particular, locally free implies
reflexive, and reflexive implies
torsion-free.

The reverse implications fail in
three dimensions. For a line
$L\subset\mathbb{P}^3$,
$$
\mathcal{I}_L^{**}=\mathcal{O}
\ne\mathcal{I}_L.
$$
Thus its ideal sheaf is torsion-free
but not reflexive. A reflexive sheaf
on a smooth surface is locally free;
on a smooth threefold it can fail to
be locally free at finitely many points.

\Needspace{10\baselineskip}
\textbf{Twists: the arithmetic of degrees.}
Twisting means tensoring with a line
bundle:
$$
F(a)=F\otimes\mathcal{O}(a),\qquad
F(a)(b)=F(a+b).
$$
It preserves exact sequences and rank.
Two useful sign checks are
$$
\begin{aligned}
\operatorname{Hom}(\mathcal{O}(-a),F)
&=H^0(F(a)),\\
\operatorname{Hom}
(\mathcal{O}(a),\mathcal{O}(b))
&=H^0(\mathcal{O}(b-a)).
\end{aligned}
$$
The second map is multiplication by
a homogeneous polynomial of degree
$b-a$. It must be zero if $b-a<0$.
This is the same degree bookkeeping
as the shifts $R[-a]$ in a graded
free resolution: its generator has
degree $a$.
For a rank-$r$ bundle, the determinant
picks up the twist $ra$, because each
of the $r$ factors in a top wedge
picks up one copy of $\mathcal{O}(a)$.

An inclusion is \emph{saturated} if
its quotient is torsion-free.
Multiplication by a linear equation
$\ell$ gives an injection
$\mathcal{O}(-1)\to\mathcal{O}$,
but its quotient is $\mathcal{O}_H$,
where $H=V(\ell)$. This quotient has
torsion on projective space, so the
inclusion is not saturated.
In contrast, the quotient in (1) is
an ideal sheaf, hence torsion-free.

\Needspace{8\baselineskip}
\textbf{Chern classes in the calculations.}
These are the same characteristic
classes used in differential geometry.
For a complex vector bundle,
$c_i(E)\in H^{2i}(X,\mathbb{Z})$.
On $\mathbb{P}^3$, all we need is
$$
H^*(\mathbb{P}^3,\mathbb{Z})
=\mathbb{Z}[h]/(h^4).
$$
Write $c_1(E)=ch$ and $c_2(E)=eh^2$.
For rank two,
$$
\begin{aligned}
c_1(E(a))&=(c+2a)h,\\
c_2(E(a))&=(e+ac+a^2)h^2.
\end{aligned}
$$
To remember the second formula, use
formal Chern roots: twisting adds
$ah$ to both roots. Expanding their
product gives the terms $e$, $ac$
and $a^2$.
For a regular section of $E(a)$ with
curve zero scheme $C$, its zero-locus
class is the top Chern class:
$$
[C]=c_2(E(a)),\qquad
\deg C=e+ac+a^2.
$$

For coherent sheaves on smooth
projective space, Chern classes are
defined using finite resolutions by
bundles. Their total Chern classes
multiply in short exact sequences.
This is why we divide the Chern
polynomials in the line and conic
calculations.
A rank-two bundle has $c_3=0$.
A rank-two reflexive sheaf can have
$c_3\ne0$: its class comes from an
alternating sum of bundles in a
resolution, not from one bundle of
rank two. Here nonzero $c_3$ is an
obstruction to local freeness.

\Needspace{9\baselineskip}
\textbf{Global Ext: assembling a middle
sheaf and detecting splitting.}
The group $\operatorname{Ext}^1(B,A)$
classifies short exact sequences
$$
0\longrightarrow A\longrightarrow E
\longrightarrow B\longrightarrow0,
$$
up to isomorphisms that are the identity
on $A$ and $B$. The zero class is the
split sequence, with $E=A\oplus B$.
Applying $\operatorname{Hom}(-,A)$
gives part of a long exact sequence:
$$
\operatorname{Hom}(E,A)
\longrightarrow\operatorname{Hom}(A,A)
\xrightarrow{\delta}
\operatorname{Ext}^1(B,A).
$$
The image $\delta(1_A)$ is the extension
class. Exactness says that it vanishes
exactly when $1_A$ extends to a map
$E\to A$. That map is a retraction of
$A\to E$, so the sequence splits.
This explains how Ext enters when
we apply Hom to a short exact sequence.

Over a ring, one can compute Ext by
resolving the first argument by
projective modules, or the second by
injective modules. These give the same
groups. For sheaves, a locally free
resolution computes a sheaf Hom
complex; global Ext also requires
its hypercohomology. In particular,
if $E$ is locally free, then
$$
\operatorname{Ext}^i(E,A)
\cong H^i(X,E^*\otimes A).
$$
Local freeness of the first argument
therefore does not force global Ext
to vanish.

\Needspace{10\baselineskip}
\textbf{The Ext-to-cohomology calculation.}
On $\mathbb{P}^3$, Serre duality says
$$
\operatorname{Ext}^1(B,A)
\cong\operatorname{Ext}^2
(A,B(-4))^\vee.
$$
The degree $1$ becomes $3-1=2$;
the arguments exchange places; and
$\omega_X=\mathcal{O}(-4)$ contributes
the twist $-4$.
Insert $B=\mathcal{I}_C(c+a)$ and
$A=\mathcal{O}(-a)$:
$$
\begin{aligned}
\operatorname{Ext}^1
(\mathcal{I}_C(c+a),\mathcal{O}(-a))
&\cong\operatorname{Ext}^2
(\mathcal{O}(-a),\mathcal{I}_C(c+a-4))^\vee
\\
&\cong H^2
(\mathcal{I}_C(c+2a-4))^\vee.
\end{aligned}
$$
The last step is useful because the
first argument is now a line bundle.
Its dual is $\mathcal{O}(a)$; tensoring
by that dual supplies the second $a$.
Thus all the twists are accounted for.

Next use
$0\to\mathcal{I}_C(m)\to\mathcal{O}(m)
\to\mathcal{O}_C(m)\to0$.
The long exact cohomology sequence
contains
$$
0\longrightarrow H^1(\mathcal{O}_C(m))
\longrightarrow H^2(\mathcal{I}_C(m))
\longrightarrow0,
$$
since $H^1(\mathcal{O}(m))$ and
$H^2(\mathcal{O}(m))$ vanish on
$\mathbb{P}^3$ for every $m$.
So the two middle groups are isomorphic.
For a line or conic this leaves a
calculation on $\mathbb{P}^1$, where
$$
\begin{aligned}
h^1(\mathcal{O}_{\mathbb{P}^1}(n))
&=h^0(\mathcal{O}_{\mathbb{P}^1}(-n-2))\\
&=\max\{0,-n-1\}.
\end{aligned}
$$
The vanishing used above concerns
\emph{intermediate cohomology of line
bundles on projective three-space}.
It says nothing of that kind for
arbitrary sheaves or for top cohomology.
Indeed $H^3(\mathcal{O}(-4))
\cong\mathbb{C}$.

\Needspace{10\baselineskip}
\textbf{Sheaf Ext: a local test.}
The sheaf $\mathcal{E}xt^i(F,G)$ has
stalks
$$
\mathcal{E}xt^i(F,G)_x
\cong\operatorname{Ext}^i_{\mathcal{O}_{X,x}}
(F_x,G_x).
$$
One computes it by applying sheaf Hom
to a locally free resolution and taking
the cohomology sheaves of that complex.
Equivalently, it is the sheafification
of the presheaf of Ext groups on open
sets. It records local information,
including where its stalks are nonzero.

If $F$ is a bundle, its stalks are free
modules, so its positive sheaf Ext
vanishes. This is different from global
Ext: for example,
$$
\operatorname{Ext}^3
(\mathcal{O},\mathcal{O}(-4))
\cong H^3(\mathcal{O}(-4))
\cong\mathbb{C},
$$
although the corresponding positive
sheaf Ext groups vanish.

\newpage
\textbf{Charles's local-freeness and
reflexivity criteria.}
Let $F$ be torsion-free on a smooth
projective threefold, and abbreviate
$$
\mathcal{E}^i=
\mathcal{E}xt^i(F,\mathcal{O}_X).
$$
Then the precise statements are:
\begin{itemize}
\item $F$ is locally free if and only
if $\mathcal{E}^1=\mathcal{E}^2=0$.
\item $F$ is reflexive if and only
if $\mathcal{E}^2=0$ and
$\mathcal{E}^1$ has finite support.
\end{itemize}
Finite support permits a nonzero stalk
at finitely many points, or none.
Thus, once we know that $F$ is reflexive,
$\mathcal{E}^1=0$ is exactly the test
for local freeness. A reflexive sheaf
may fail this test at isolated points,
but not along a curve.

For the sheaves in the line-and-conic
exercise, the ideal resolutions and
the extension by a line bundle give
local free resolutions of length at
most one. Hence $\mathcal{E}^2=0$
is automatic there; finite support
of $\mathcal{E}^1$ tests reflexivity.
For a general torsion-free sheaf the
condition on $\mathcal{E}^2$ must also
be checked.

The general support criterion for
reflexivity is \cite{hl}
$$
\operatorname{codim}_X\operatorname{Supp}
\mathcal{E}xt^i(F,\mathcal{O}_X)
\geq i+2\qquad(i>0),
$$
with $F$ torsion-free. For $i=1$ the
support must have codimension at least
three, so on a projective threefold
it is finite. For $i=2$ it would need
codimension at least four, so it must
be empty.
Locally, Auslander--Buchsbaum says
$$
\operatorname{pd}F_x+
\operatorname{depth}F_x
=\dim\mathcal{O}_{X,x}.
$$
Projective dimension is the length of
a shortest free resolution of the
stalk. Torsion-freeness gives positive
depth at points of positive codimension,
so this length is at most two here.
The highest nonzero local Ext detects
this length. This explains why only
$\mathcal{E}^1$ and $\mathcal{E}^2$ are
needed in the two tests.

Finally, the ideal sheaf of a point
shows why checking only first Ext
can give the wrong answer:
$$
\mathcal{E}xt^1(\mathcal{I}_p,
\mathcal{O})=0,\qquad
\mathcal{E}xt^2(\mathcal{I}_p,
\mathcal{O})\cong\mathcal{O}_p.
$$
Nevertheless $\mathcal{I}_p$ is neither
locally free nor reflexive.

\begin{thebibliography}{3}
\raggedright
\bibitem{rao}
R.~Hartshorne,
\emph{On Rao's Theorems and the
Lazarsfeld--Rao Property}, 2003.
\href{https://arxiv.org/abs/math/0302078}
{arXiv:math/0302078}.

\bibitem{hl}
D.~Huybrechts and M.~Lehn,
\emph{The Geometry of Moduli Spaces
of Sheaves}, Section~1.1,
especially Propositions~1.1.6 and~1.1.10.

\bibitem{stacks}
The Stacks Project,
\href{https://stacks.math.columbia.edu/tag/0AV0}
{Tag 0AV0} and
\href{https://stacks.math.columbia.edu/tag/0AVT}
{Tag 0AVT}, on reflexive modules and
sheaves.
\end{thebibliography}
\end{document}
